\section{Анализ на сигурността}
\label{sec:security}

\subsection{Модел на заплахите}

Нападателят се разглежда като произволен клиент в мрежата. Той може да отваря връзки,
да изпраща каквото пожелае, когато пожелае, и да не изпраща нищо. Нападател, който е
между клиента и сървъра и може да променя трафика, се разглежда отделно в
подраздел~\ref{sec:downgrade}, защото там проблемът не е в кода, а в самата идея.

Активът е сървърът: неговите нишки, памет и процесорно време, както и коректността на
решението кой протокол се говори. Целта на анализа е за всяка заплаха да посочи кое
място в кода отговаря на нея и кой тест го проверява.

\subsection{Заплахи и отговори}

\begin{longtable}{@{}L{4.6cm}L{5.8cm}L{5.4cm}@{}}
\caption{Заплахи и отговорът на \code{octet}}
\label{tab:threats}\\
\toprule
Заплаха & Отговор в кода & Проверено от \\
\midrule
\endfirsthead
\toprule
Заплаха & Отговор в кода & Проверено от \\
\midrule
\endhead
Клиент се свързва и не изпраща нищо & общ краен срок от приемането; връзката се
  затваря при изтичането му & тест: затваряне при срока, докато други клиенти се
  обслужват \\
Клиент изпраща 0x16 и спира & същият срок покрива и ръкостискането & тест: затваряне
  при същия срок \\
Клиент изпраща по един октет точно преди изтичане & срокът е абсолютен и не се
  удължава от получени данни & по построение (листинг~\ref{lst:deadline}) \\
Много мълчащи връзки наведнъж & горна граница на едновременните връзки; над нея новите
  се затварят веднага & тест: третата връзка при граница 2 \\
0x16, последван от безсмислени данни & разпознаването насочва към TLS, OpenSSL
  отхвърля записа & тест: няма некриптиран отговор \\
Произволни двоични данни или малки букви & „непознато“, затваряне без отговор & тестове
  за двоични данни и за \code{get} \\
Предупреждение (0x15) като първи запис & „непознато“, не се подава на OpenSSL & тест \\
Главна буква, която не е метод & анализаторът приема формата, сървърът отговаря 501
  \cite[раздел 15.6.2]{rfc9110} & тест с метод \code{ZZZ} \\
Неправилен ред на заявката & отговор 400 & тест \\
Огромна заглавна част на заявката & граница 8192 октета, отговор 431 \cite{rfc6585}
  & тест; виж подраздел~\ref{sec:bug} \\
Клиент без общ приложен протокол & фатално предупреждение 120 &
  тест по октетите в мрежата \\
Клиент със стара версия & минимална версия на контекста & тест с клиент до TLS 1.2
  срещу сървър с минимум 1.3 \\
TLS клиент към сървър без сертификат & връзката се затваря; некриптиран отговор
  не би бил разбран & тест \\
Сървър с чуждо име & проверката е у клиента; сървърът не може да я направи вместо него
  & тест в паметта: грешка за несъвпадащо име \\
\bottomrule
\end{longtable}

Няколко реда от таблицата изискват обяснение.

\textbf{Защо непознатото се затваря без отговор.} Връзка, която не започва нито с TLS
запис, нито с HTTP метод, не говори протокол, който сървърът разбира. Всеки отговор би
бил на език, който другата страна не е използвала. Освен това отговорът дава на
сканиращия нападател информация за сървъра. Мълчаливото затваряне не дава нищо.

\textbf{Защо разпознаването не проверява повече от един октет.} Строг разпознавач би
могъл да провери още версията на записа, границата на дължината и типа на съобщението
за ръкостискане. Това изисква до шест октета и следователно чакане, докато пристигнат.
Тези проверки обаче OpenSSL прави така или иначе и ги прави по-добре: той знае всички
правила на протокола, а разпознавачът би знаел само някои от тях. Разпознаването в
\code{octet} насочва, а не валидира. Тестът „0x16, последван от безсмислени данни“
проверява точно това: насочването е към TLS, но връзката не получава некриптиран
отговор, защото OpenSSL отхвърля записа.

\textbf{Защо проверката на името е у клиента.} Сертификатът доказва, че сървърът
притежава ключ за дадено име. Кое име е очаквал клиентът, знае само клиентът. Затова
тестът със сървър с чуждо име се проваля в клиента, с грешка за несъвпадащо име, а не в
сървъра.

\subsection{Отклонение от препоръка в HTTP/1.1}

RFC 9112 препоръчва сървърът да игнорира поне един празен ред, получен преди реда на
заявката \cite[раздел 2.2]{rfc9112}. Препоръката е формулирана като SHOULD, тоест не е
задължителна. \code{octet} не я следва: връзка, която започва с CR (0x0D), не започва с
главна буква и се затваря като непозната. Това е съзнателен избор. Ако се приемат
празни редове преди заявката, разпознаването вече не може да реши по първия октет и
трябва да чака, докато празните редове свършат. Клиент, който изпраща празен ред преди
заявката си, няма да бъде обслужен от \code{octet}.

\subsection{Грешката, която тестовете откриха}
\label{sec:bug}

Първата версия на анализатора на HTTP проверяваше границата от 8192 октета само докато
не намери празния ред, който завършва заглавната част. Ако едно четене донесе
едновременно октетите над границата и празния ред, проверката за празния ред се
изпълняваше първа и заявката се приемаше. Тестът на ниво функция не откри грешката,
защото подава заявката по един октет. Интеграционният тест я откри: заявка с поле от 9000
октета получи отговор 200 вместо 431. Поправката проверява позицията на празния ред
спрямо границата, а нов тест на ниво функция подава заявката на големи части, както я
доставя мрежата.

\subsection{Атаки с понижаване на защитата}
\label{sec:downgrade}

Това е най-сериозното ограничение на темата и то не е в кода. RFC 9325 казва директно,
че съчетаването на незащитена и защитена с TLS комуникация отваря пътя за SSL Stripping и
подобни атаки, защото началото на комуникацията не е защитено от промяна и нападателят
може да го манипулира така, че комуникацията да остане некриптирана
\cite[раздел 3.2]{rfc9325}. SSL Stripping е описан в RFC 7457 \cite[раздел 2.1]{rfc7457}:
нападател между клиента и сървъра премахва преминаването към TLS от некриптирания HTTP
трафик.

Сървър, който обслужва и TLS, и некриптиран HTTP на един порт, е точно такова
съчетание. Разпознаването по първия октет не създава уязвимостта: тя съществува и при
два отделни порта, ако на порт 80 се обслужва съдържание. Но разпознаването прави
некриптирания достъп до същата услуга възможен по подразбиране.

Стандартната защита е HTTP Strict Transport Security \cite{rfc6797}. Сървърът изпраща
по TLS заглавно поле, с което казва на браузъра за определен период да се свързва с него
само по TLS. RFC 9325 изисква уеб сървърите да го поддържат \cite[раздел 3.2]{rfc9325}.
\code{octet} в сегашния си вид не изпраща това заглавно поле и отговаря с едно и също
съдържание и по двата транспорта. За производствена употреба правилната схема е
следната: некриптираният път само пренасочва към адреса по TLS, а отговорите по TLS
носят заглавното поле за HSTS. Тогава ползата от един порт остава (един сокет, една
конфигурация), а некриптираният път не обслужва съдържание.

\subsection{Какво не е проверено}

Сървърът не е подлаган на автоматизирано подаване на случайни входни данни (fuzzing),
нито на одит от трета страна. Кодът на OpenSSL се приема за коректен. Нишка за всяка
връзка е проста схема, но тя ограничава броя на едновременните връзки до зададената
граница. Сървър, който трябва да издържа на големи атаки за отказ от обслужване, би
използвал асинхронен вход и изход. RFC 7457 отбелязва, че TLS сам по себе си е вектор за
такива атаки, защото клиентът с малко усилие кара сървъра да извърши значителна
изчислителна работа \cite[раздел 2.13]{rfc7457}. Това важи и за \code{octet}.
